\section{仿射簇}\label{sec:affine-varieties}

并非 $\mathbb{A}^{n}$ 的每个子集都是某族多项式的零点集. 例如, 设 $X\subseteq\mathbb{A}^{1}_{\mathbb{R}}$ 为全体非零实数的集合:
\[
  X=\{a\in\mathbb{R}\mid a\ne0\}\subsetneq\mathbb{A}^{1}_{\mathbb{R}}.
\]
我们来说明 $X$ 不能是任何多项式集合的零点集. 反设对某个集合 $S\subseteq\mathbb{R}[x]$, 有 $X=\mathcal{V}(S)$. 那么 $S$ 必须至少包含一个非零元素, 因为如果 $S$ 是 $\varnothing$ 或 $\{0\}$, 它的零点集就是整个 $\mathbb{A}^{1}_{\mathbb{R}}$. 取 $S$ 中任一个非零元素 $f$. 由于 $X=\mathcal{V}(S)$, 对所有 $a\in\mathbb{R}\setminus\{0\}$ 都有 $f(a)=0$. 这意味着 $f\in\mathbb{R}[x]$ 是一个有无穷多个零点的非零多项式, 与 Corollary~\ref{cor:polynomial-finite-zeros} 矛盾.

鉴于这类例子, 我们给仿射空间中能定义为一个多项式集合的零点集的那些特殊子集起一个名字.

\begin{Definition}{仿射簇}\leavevmode\label{def:affine-variety}
  若对某个多项式集合 $S\subseteq K[x_{1},\ldots,x_{n}]$, 有 $X=\mathcal{V}(S)$, 则称子集 $X\subseteq\mathbb{A}^{n}$ 为\textbf{仿射簇}.
\end{Definition}

我们已经在第~\ref{sec:vanishing-set-operator}~节给出的例子中见过若干仿射簇. 此外, 对第~\ref{sec:vanishing-set-operator}~节中的每个仿射簇, 我们都给出了一个以它为零点集的多项式集合. 另外两个例子可能比前面的任何例子都更基本, 却同样重要: 空集和整个仿射空间.

\begin{Example}{空集和仿射空间都是仿射簇}\leavevmode\label{ex:empty-and-full-affine-varieties}
  常数多项式 $1\in K[x_{1},\ldots,x_{n}]$ 在 $\mathbb{A}^{n}$ 的任一点都不为零, 所以
  \[
    \mathcal{V}(1)=\varnothing\subseteq\mathbb{A}^{n}.
  \]
  零多项式 $0\in K[x_{1},\ldots,x_{n}]$ 在 $\mathbb{A}^{n}$ 的每一点都为零, 所以
  \[
    \mathcal{V}(0)=\mathbb{A}^{n}.
  \]
  因而 $\varnothing$ 和 $\mathbb{A}^{n}$ 都是仿射簇.
\end{Example}

只要写出一些多项式集合, 再考察它们的零点集, 就很容易列出无穷多个仿射簇的例子, 不过这种做法并不能说明太多问题. 更微妙的任务是弄清仿射簇在仿射空间的子集中有什么特别之处. 换句话说, 哪些仿射空间的子集不是仿射簇? 我们已经见过一个例子: $\mathbb{R}\setminus\{0\}$ 不是 $\mathbb{A}^{1}_{\mathbb{R}}$ 中的仿射簇. 事实上, $\mathbb{A}^{1}$ 中仅有的仿射簇是有限点集 (包括 $\varnothing$) 和整个 $\mathbb{A}^{1}$ (Exercise~\ref{ex:affine-varieties-on-the-line}).

另一方面, $\mathbb{A}^{2}$ 中的无限真子集完全可以是仿射簇, $\mathbb{A}^{2}_{\mathbb{R}}$ 中的抛物线就是一例. 但 $\mathbb{A}^{2}_{\mathbb{R}}$ 的真子集要成为仿射簇, 就不可能有非空内部. 下面的例子说明这一现象.

\begin{Example}{实心正方形不是仿射簇}\leavevmode\label{ex:solid-square-not-affine-variety}
  设 $X$ 是 $\mathbb{A}^{2}_{\mathbb{R}}$ 中如下定义的实心正方形:
  \[
    X=\{(a,b)\in\mathbb{A}^{2}_{\mathbb{R}}\mid-1\leq a\leq1
      \text{ 且 }-1\leq b\leq1\}.
  \]
  反设对某个集合 $S\subseteq\mathbb{R}[x,y]$, 有 $X=\mathcal{V}(S)$. 任取 $f\in S$, 我们来说明 $f$ 是零多项式. 写成
  \[
    f=\sum_{i=0}^{d}f_{i}y^{i},
  \]
  其中对每个 $i$ 都有 $f_{i}\in\mathbb{R}[x]$. 由于 $f\in S$, 它在 $X$ 的所有点上都为零. 换句话说, 对任意 $a\in[-1,1]$, 一元多项式
  \[
    f(a,y)=\sum_{i=0}^{d}f_{i}(a)y^{i}\in\mathbb{R}[y]
  \]
  在所有 $b\in[-1,1]$ 处都为零. 由于非零一元多项式只能有有限多个零点, 可知 $f(a,y)$ 必为零多项式, 从而对所有 $i$ 都有 $f_{i}(a)=0$. 换句话说, 我们已经说明一元多项式 $f_{i}\in\mathbb{R}[x]$ 在所有 $a\in[-1,1]$ 处都为零. 再次利用非零一元多项式只有有限多个零点这一事实, 可知对每个 $i$, $f_{i}$ 都是零多项式, 所以 $f$ 是零多项式. 这一论证表明, $S$ 中可能有的多项式只有零多项式, 从而 $X=\mathcal{V}(S)=\mathbb{A}^{2}_{\mathbb{R}}$, 矛盾.

  学过拓扑学的学生不妨证明一个更一般的结论: 赋予 $\mathbb{A}^{n}_{\mathbb{R}}$ 欧氏拓扑后, $\mathbb{A}^{n}_{\mathbb{R}}$ 中唯一具有非空拓扑内部的仿射簇是整个 $\mathbb{A}^{n}_{\mathbb{R}}$ (Exercise~\ref{ex:affine-variety-empty-interior}).
\end{Example}

\begin{Example}{$\mathbb{A}^{2}_{\mathbb{R}}$ 中 $e^{x}$ 的图像不是仿射簇}\leavevmode\label{ex:exponential-graph-not-affine-variety}
  令 $X$ 为实数上的指数函数的图像:
  \begin{equation}\label{eq:exponential-graph}
    X=\{(a,b)\in\mathbb{A}^{2}_{\mathbb{R}}\mid b=e^{a}\}\subseteq\mathbb{A}^{2}_{\mathbb{R}}.
  \end{equation}
  Exercise~\ref{ex:prove-exponential-graph-not-affine-variety} 给出了证明 $X$ 不是仿射簇的严谨思路. 直观上, $X$ 不是仿射簇似乎是合理的: 它由表达式 $y-e^{x}$ 为零来定义, 而这个表达式不是 $x$ 和 $y$ 的多项式. 不过, 采用这种推理时要谨慎. 例如, 表达式 $\sin^{2}(x)+\cos^{2}(x)+y$ 也不是多项式, 但集合
  \[
    \{(a,b)\in\mathbb{A}^{2}_{\mathbb{R}}\mid\sin^{2}(a)+\cos^{2}(a)+b=0\}
    \subseteq\mathbb{A}^{2}_{\mathbb{R}}
  \]
  与集合
  \[
    \{(a,b)\in\mathbb{A}^{2}_{\mathbb{R}}\mid1+b=0\}=\mathcal{V}(1+y)
    \subseteq\mathbb{A}^{2}_{\mathbb{R}}
  \]
  相同, 所以它是一个仿射簇.
\end{Example}

不同的集合 $S$ 的 $\mathcal{V}(S)$ 完全可能给出同一个仿射簇. 例如,
\[
  \mathcal{V}(x,y)=\mathcal{V}(x+y,x-y)=\{(0,0)\}\subseteq\mathbb{A}^{2},
\]
读者不难验证这一点. 事实上, 可以用 $S$ 生成的整个理想 $\langle S\rangle\subseteq K[x_{1},\ldots,x_{n}]$ 替换 $S$, 而不改变其零点集.

\begin{Proposition}{仿射簇由理想定义}\leavevmode\label{prop:affine-varieties-defined-by-ideals}
  若 $S\subseteq K[x_{1},\ldots,x_{n}]$ 是一个多项式集合, 则
  \[
    \mathcal{V}(S)=\mathcal{V}(\langle S\rangle).
  \]
\end{Proposition}

\begin{Proof}
  见 Exercise~\ref{ex:prove-ideals-define-affine-varieties}.
\end{Proof}

为了概括本节的讨论, 我们把 $\mathcal{V}$ 算子看成函数
\[
  \mathcal{V}:\{K[x_{1},\ldots,x_{n}]\text{ 的子集}\}
    \longrightarrow\{\mathbb{A}^{n}\text{ 的子集}\}.
\]
这个函数不是满射, 因为 $\mathbb{A}^{n}$ 的子集不全是仿射簇. 但如果把陪域限制为仿射簇, 则根据定义, 它就成为满射:
\[
  \mathcal{V}:\{K[x_{1},\ldots,x_{n}]\text{ 的子集}\}
    \longrightarrow\{\mathbb{A}^{n}\text{ 中的仿射簇}\}.
\]
$\mathcal{V}$ 算子也不是单射, 因为 $K[x_{1},\ldots,x_{n}]$ 的不同子集可以定义同一个仿射簇. 为了使它成为双射, 第一步是利用 Proposition~\ref{prop:affine-varieties-defined-by-ideals}: 可以把定义域限制为理想, 同时保持函数是满射:
\[
  \mathcal{V}:\{K[x_{1},\ldots,x_{n}]\text{ 的理想}\}
    \longrightarrow\{\mathbb{A}^{n}\text{ 中的仿射簇}\}.
\]
事实上, 这个算子仍不是单射 (见 Exercise~\ref{ex:distinct-ideals-same-variety}). 要最终得到一个双射的 $\mathcal{V}$ 算子, 还需要进一步限制定义域, 并对基域 $K$ 作一个假设.

\BookExerciseHeading
\begin{BookExercises}
  \item\label{ex:affine-varieties-on-the-line} 证明 $\mathbb{A}^{1}$ 中仅有的仿射簇是有限点集和整个 $\mathbb{A}^{1}$.
  \item\label{ex:prove-ideals-define-affine-varieties} 证明 Proposition~\ref{prop:affine-varieties-defined-by-ideals}.
  \item\label{ex:distinct-ideals-same-variety} 举出不同的理想 $I,J\in K[x]$, 使得 $\mathcal{V}(I)=\mathcal{V}(J)$.
  \item\label{ex:twisted-cubic-is-affine-variety} 求一个多项式集合 $S$, 使得 $X=\mathcal{V}(S)$, 由此证明集合
    \[
      X=\{(a,a^{2},a^{3})\mid a\in K\}\subseteq\mathbb{A}^{3}
    \]
    是一个仿射簇. (簇 $X$ 称为仿射扭三次曲线.)
  \item\label{ex:containment-in-affine-plane} 设 $K$ 是无限域. $\mathbb{A}^{3}_{K}$ 中的一个平面是能定义为某个非常数线性多项式的零点集的仿射簇 $P$:
    \[
      P=\mathcal{V}(Ax+By+Cz+D),
    \]
    其中 $A,B,C,D\in K$, 且 $A,B,C$ 不全为零. 证明下列结论.
    \begin{enumerate}[label=(\alph*)]
      \item 对任意次数至多为二的多项式 $f_{1},f_{2},f_{3}\in K[x]$, 集合
        \[
          X=\{(f_{1}(a),f_{2}(a),f_{3}(a))\mid a\in K\}\subseteq\mathbb{A}^{3}
        \]
        至少包含在一个平面内. (提示: 对 $A,B,C,D$ 建立线性方程, 并说明至少存在一个解.)
      \item Exercise~\ref{ex:twisted-cubic-is-affine-variety} 中的仿射扭三次曲线不包含在任何平面内. (因此, 把 (a) 中的\textquotedbl{}二\textquotedbl{}换成\textquotedbl{}三\textquotedbl{}后, 结论不成立.)
    \end{enumerate}
  \item\label{ex:parameterized-set-not-affine-variety} 设 $K$ 是无限域. 证明集合
    \[
      X=\{(a,ab)\mid a,b\in K\}\subseteq\mathbb{A}^{2}
    \]
    不是仿射簇.
  \item\label{ex:prove-exponential-graph-not-affine-variety} 证明表达式 $y-e^{x}$ 在 $\mathbb{A}^{2}_{\mathbb{R}}$ 中的零点集 $X$ 不是仿射簇, 可以采用以下策略.
    \begin{enumerate}[label=(\alph*)]
      \item 反设对某个 $S\subseteq\mathbb{R}[x,y]$, 有 $X=\mathcal{V}(S)$. 说明为什么存在非零多项式 $f\in\mathbb{R}[x,y]$, 使得对所有 $a\in\mathbb{R}$ 都有 $f(a,e^{a})=0$.
      \item 对上述任意多项式 $f$, 写成
        \[
          f(x,y)=p_{0}(x)+p_{1}(x)y+p_{2}(x)y^{2}+\cdots+p_{d}(x)y^{d},
        \]
        其中 $p_{d}\ne0\in\mathbb{R}[x]$. 证明对所有 $a\in\mathbb{R}$ 都有
        \[
          \frac{p_{0}(a)}{e^{da}}+\frac{p_{1}(a)}{e^{(d-1)a}}+
          \frac{p_{2}(a)}{e^{(d-2)a}}+\cdots+
          \frac{p_{d-1}(a)}{e^{a}}+p_{d}(a)=0.
        \]
      \item 令 $a\to\infty$, 取上式的极限, 得到
        \[
          \lim_{a\to\infty}p_{d}(a)=0.
        \]
        说明对任意非零的 $g\in\mathbb{R}[x]$, 有 $\lim g(a)\ne0$, 从而得出矛盾.
    \end{enumerate}
  \item\label{ex:affine-variety-empty-interior} (供学过一些拓扑学的学生.) 把 $\mathbb{A}^{n}_{\mathbb{R}}=\mathbb{R}^{n}$ 看作带欧氏拓扑的拓扑空间. 设 $X\subseteq\mathbb{A}^{n}_{\mathbb{R}}$ 是一个不同于 $\mathbb{A}^{n}_{\mathbb{R}}$ 本身的仿射簇. 证明: 作为 $\mathbb{A}^{n}_{\mathbb{R}}$ 的拓扑子空间, $X$ 的内部为空. (用类似的证明, 把 $\mathbb{R}$ 换成 $\mathbb{C}$ 后, 同一结果也成立.)
\end{BookExercises}
